% Archivo generado automáticamente por simplex.c
\documentclass[12pt]{article}
\usepackage[utf8]{inputenc}
\usepackage{graphicx}
\usepackage{array,booktabs}
\usepackage[table]{xcolor}
\usepackage{longtable}
\usepackage{geometry}
\usepackage{pdflscape}
\usepackage{tikz}
\geometry{margin=0.8in}
\begin{document}
\begin{center}
{\large Instituto Tecnológico de Costa Rica\\[1cm]
\includegraphics[width=0.4\textwidth]{TEC.png}\\[2cm]
{\LARGE \textbf{Proyecto 4: SIMPLEX}}\\[2cm]
{\large Investigación de Operaciones\\[2cm]
{\large Profesor: }\\[1cm]
{\large Francisco Jose Torres Roja}\\[2cm]
{\large Integrantes: }\\[1cm]
{\large Jose Pablo Fernandez Jimenez - 2023117752}\\[1cm]
{\large Diego Durán Rodríguez - 2022437509}\\[2cm]
{\large Segundo semestre 2025\\[1cm]
\end{center}
\newpage
\section*{Algoritmo SIMPLEX}
El \textbf{algoritmo SIMPLEX} es un método utilizado para resolver problemas de programación lineal. 

El objetivo es minimizar o maximizar el costo total esperado. 

\textbf{Variantes del problema:}
\begin{itemize}
  \item \textit{Horizonte finito vs. infinito:} El análisis puede hacerse en un período limitado de tiempo o indefinido.
  \item \textit{Determinístico vs. estocástico:} En la versión determinística se conocen los costos y valores de reventa; en la estocástica, se modelan como variables aleatorias.
  \item \textit{Reemplazo individual vs. múltiple:} Puede plantearse para un único equipo o para varios equipos en paralelo.
\end{itemize}

\bigskip
\section*{Problema original}
El problema original se puede formular como un problema de programación lineal, donde se busca optimizar una función objetivo sujeta a ciertas restricciones.

\bigskip
\section*{Tabla inicial}
\newpage
\section*{Tabla final}
\section*{Soluciones}
\begin{thebibliography}{9}
\bibitem{meyer1971} Meyer, R. A. (1971). Equipment replacement under uncertainty. \textit{Management Science, 17}(11), 750--758. https://doi.org/10.1287/mnsc.17.11.750

\bibitem{tan2010} Tan, C., \\& Hartman, J. (2010). Equipment replacement analysis with an uncertain finite horizon. Disponible en: https://econpapers.repec.org/article/tafuiiexx/v\_3a42\_3ay\_3a2010\_3ai\_3a5\_3ap\_3a342-353.htm

\end{thebibliography}
\end{document}
